\subsubsection*{Solution}
\paragraph*{Step-by-Step Derivation}

We use the phasor convention

\[
\mathbf E(\mathbf r,t)=\operatorname{Re}\!\left[\mathbf E(\mathbf r)e^{-i\omega t}\right],
\]

and allow complex scalar polarizabilities

\[
\alpha_j=|\alpha_j|e^{i\psi_j}.
\]

They are treated as fixed, physical (radiatively dressed) polarizabilities at
the laser frequency.

Define the phase of their product by

\[
\psi_\alpha\equiv\psi_1+\psi_2=\arg(\alpha_1\alpha_2).
\]

\subparagraph*{1. Gaussian-beam phase near the focal plane}

On the propagation axis, the field of tweezer $j$ is

\[
\mathbf E_j(z_j) = \hat{\mathbf e}\, \frac{E_j}{\sqrt{1+(z_j/z_R)^2}} \exp\!\left[ i\left(kz_j-\arctan\frac{z_j}{z_R}+\phi_j\right) \right].
\]

Because $|z_j|\ll z_R$,

\[
\frac{1}{\sqrt{1+(z_j/z_R)^2}} =1+O\!\left(\frac{z_j^2}{z_R^2}\right),
\]

and

\[
\arctan\frac{z_j}{z_R} = \frac{z_j}{z_R} +O\!\left(\frac{z_j^3}{z_R^3}\right).
\]

Thus, to linear order,

\[
\mathbf E_j(z_j) \simeq \hat{\mathbf e}\,E_j e^{i(qz_j+\phi_j)}, \qquad q\equiv k-\frac{1}{z_R}.
\]

Consequently, the local relative phase is

\[
\Delta\phi = \phi_1-\phi_2+q(z_1-z_2).
\]

This longitudinal phase dependence produces the leading $z$-direction
coupling derived below.

\subparagraph*{2. Far-field dipole Green tensor}

For complex polarizabilities, the two second-order contributions to the
time-averaged optical-binding force cannot generally be combined into the
single gradient used for real polarizabilities.  In component form they are

\[
\begin{aligned}
(F_j^{\mathrm{bind}})_a
=\frac{1}{2\varepsilon_0}\operatorname{Re}\!\Big[{}
&\alpha_j^*\alpha_\ell\,
\mathbf E_j^* \cdot
\partial_{j,a}\!\left(\mathbf G_{j\ell}\mathbf E_\ell\right)\\
&+\alpha_j^*\alpha_\ell^*\,
(\mathbf G_{j\ell}\mathbf E_\ell)^* \cdot
\partial_{j,a}\mathbf E_j
\Big],
\qquad \ell\neq j,
\end{aligned}
\]

where

\[
\mathbf G_{j\ell}\equiv
\mathbf G(\mathbf r_j-\mathbf r_\ell).
\]

The first term is the force of the field scattered by particle $\ell$ on the
dipole directly induced in particle $j$.  The second is the trapping field
acting on the correction to the dipole of particle $j$ induced by that
scattered field.  When the polarizabilities are real, their sum is the
gradient of the scalar introduced below.

To see what changes when the polarizabilities are complex, introduce the
auxiliary scalar

\[
\mathcal S_j\equiv
\frac{1}{2\varepsilon_0}\operatorname{Re}\!\left[
\alpha_j\alpha_\ell\,
\mathbf E_j^*\cdot\mathbf G_{j\ell}\mathbf E_\ell
\right].
\]

Before making the leading longitudinal approximation, the force obeys

\[
\begin{aligned}
(F_j^{\mathrm{bind}})_a
={}&\partial_{j,a}\mathcal S_j\\
&+\frac{1}{2\varepsilon_0}\operatorname{Re}\!\left[
(\alpha_j^*-\alpha_j)\alpha_\ell\,
\mathbf E_j^*\cdot
(\partial_{j,a}\mathbf G_{j\ell})\mathbf E_\ell
\right].
\end{aligned}
\]

Thus $\mathcal S_j$ is an exact force generator only for real $\alpha_j$ (or
when the Green-tensor-gradient term vanishes).  For complex polarizabilities,
it generates only the dominant longitudinal phase-gradient contribution used
below.

For $kd\gg1$, the free-space Green tensor becomes

\[
\mathbf G(\mathbf d) \simeq \frac{k^2e^{ikd}}{4\pi d} \left(\mathbf I-\hat{\mathbf d}\hat{\mathbf d}\right).
\]

At equilibrium,

\[
\mathbf d_0=(d_0,0,0).
\]

The polarization is perpendicular to $\mathbf d_0$, so

\[
\hat{\mathbf e}\cdot\hat{\mathbf d}_0=0,
\]

and therefore

\[
\hat{\mathbf e}\cdot \left(\mathbf I-\hat{\mathbf d}_0\hat{\mathbf d}_0\right) \hat{\mathbf e} =1.
\]

For the leading longitudinal term retained below, define

\[
C\equiv \frac{|\alpha_1\alpha_2|E_1E_2k^2}{8\pi\varepsilon_0}.
\]

In the present geometry, the two auxiliary scalars reduce at leading
far-field order to

\[
\mathcal S_1 = \frac{C}{d}
\cos(kd-\Delta\phi+\psi_\alpha),
\]

and

\[
\mathcal S_2 = \frac{C}{d}
\cos(kd+\Delta\phi+\psi_\alpha).
\]

The different signs of $\Delta\phi$ arise because the two interference
products are respectively $\mathbf E_1^*\cdot\mathbf E_2$ and
$\mathbf E_2^*\cdot\mathbf E_1$.  The additional phase $\psi_\alpha$ comes
from the complex product $\alpha_1\alpha_2$.  The two auxiliary scalars merely
organize the leading forces compactly; they are not a single interaction
potential for the full complex-polarizability force.

\subparagraph*{3. Expansion about the transverse equilibrium}

Let

\[
\delta z\equiv z_1-z_2.
\]

Since the equilibrium separation is along $x$,

\[
d=\sqrt{d_0^2+\delta z^2} =d_0+O(\delta z^2).
\]

Although $d-d_0=O(\!\delta z^2)$, differentiating the $d$-dependent Green
tensor produces a force linear in $\delta z$.  Its spring contribution scales
as $k^3/d_0^2$, whereas the retained phase-gradient contribution scales as
$k^2q^2/d_0$.  Their ratio is therefore

\[
O\!\left(\frac{k}{q^2d_0}\right)
\simeq O\!\left(\frac{1}{kd_0}\right),
\]

where the last estimate uses $q\simeq k$.  We neglect this term consistently
at leading far-field order.  This estimate is not uniform at a zero of the
relevant leading directional stiffness: there the displayed leading coupling
vanishes, and the nominally subleading term determines the residual coupling.

Let

\[
\Delta\phi_0\equiv\phi_1-\phi_2, \qquad \Delta\phi=\Delta\phi_0+q\delta z.
\]

Keeping only the leading phase-gradient contribution, the force on particle
1 is

\[
F_{1z}^{\mathrm{bind}} \simeq \frac{Cq}{d_0}
\sin\!\left(kd_0+\psi_\alpha-\Delta\phi_0-q\delta z\right).
\]

Using

\[
\sin(\theta-q\delta z) = \sin\theta-q\delta z\cos\theta+O(\delta z^2),
\]

we obtain

\[
F_{1z}^{\mathrm{bind}} \simeq F_{1z}^{(0)}
- \frac{Cq^2}{d_0}
\cos(kd_0+\psi_\alpha-\Delta\phi_0)(z_1-z_2),
\]

where

\[
F_{1z}^{(0)} = \frac{Cq}{d_0}
\sin(kd_0+\psi_\alpha-\Delta\phi_0)
\]

is the leading constant binding force, already balanced in the stated static
equilibrium.

Similarly,

\[
F_{2z}^{\mathrm{bind}} \simeq \frac{Cq}{d_0}
\sin\!\left(kd_0+\psi_\alpha+\Delta\phi_0+q\delta z\right),
\]

so that

\[
F_{2z}^{\mathrm{bind}} \simeq F_{2z}^{(0)}
+ \frac{Cq^2}{d_0}
\cos(kd_0+\psi_\alpha+\Delta\phi_0)(z_1-z_2).
\]

Because $z_1$ and $z_2$ are measured from the actual static equilibrium, the
constant terms are canceled by the static trap-force balance and do not enter
the linearized equations of motion.

\subparagraph*{4. Identification of $k_1$ and $k_2$}

Comparing with the stated equations gives

\[
k_1+k_2 = \frac{Cq^2}{d_0}
\cos(kd_0+\psi_\alpha-\Delta\phi_0),
\]

and

\[
k_1-k_2 = \frac{Cq^2}{d_0}
\cos(kd_0+\psi_\alpha+\Delta\phi_0).
\]

Adding these equations and using

\[
\cos(A-B)+\cos(A+B)=2\cos A\cos B,
\]

gives

\[
k_1 = \frac{Cq^2}{d_0}
\cos(kd_0+\psi_\alpha)\cos(\Delta\phi_0).
\]

Subtracting them and using

\[
\cos(A-B)-\cos(A+B)=2\sin A\sin B,
\]

gives

\[
k_2 = \frac{Cq^2}{d_0}
\sin(kd_0+\psi_\alpha)\sin(\Delta\phi_0).
\]

Substituting $C$ and $q$,

\[
k_1 = \frac{|\alpha_1\alpha_2|E_1E_2k^2}
{8\pi\varepsilon_0d_0}
\left(k-\frac{1}{z_R}\right)^2
\cos(kd_0+\psi_\alpha)\cos(\phi_1-\phi_2),
\]

\[
k_2 = \frac{|\alpha_1\alpha_2|E_1E_2k^2}
{8\pi\varepsilon_0d_0}
\left(k-\frac{1}{z_R}\right)^2
\sin(kd_0+\psi_\alpha)\sin(\phi_1-\phi_2).
\]

Here $k_1$ is the reciprocal part of the reduced linearized stiffness and, by
itself, derives from a quadratic coupling potential.  The coefficient $k_2$
is the nonreciprocal part.  This decomposition does not make the complete
driven, lossy optical system conservative.  With the arrow denoting source
particle to target particle,

\[
k_{2\to1}=k_1+k_2, \qquad k_{1\to2}=k_1-k_2,
\]

and these two directional stiffnesses are unequal in general.

For a paraxial Gaussian beam, $kz_R\gg1$, so

\[
k-\frac{1}{z_R}\simeq k.
\]

The leading-order far-field expressions are therefore

\[
k_1 \simeq \frac{|\alpha_1\alpha_2|E_1E_2k^4}
{8\pi\varepsilon_0d_0}
\cos(kd_0+\psi_\alpha)\cos(\phi_1-\phi_2),
\]

\[
k_2 \simeq \frac{|\alpha_1\alpha_2|E_1E_2k^4}
{8\pi\varepsilon_0d_0}
\sin(kd_0+\psi_\alpha)\sin(\phi_1-\phi_2).
\]

The prefactor has units of $\mathrm{N\,m^{-1}}$, as required for a spring
constant.  For positive real polarizabilities, $\psi_\alpha=0$ and
$|\alpha_1\alpha_2|=\alpha_1\alpha_2$.
